\section{The two methods: precise transfers and repairs}
\label{sec:methods}

\subsection*{MT-A01. What the adopted Dirichlet theorem says adelically}
The supplied declarations \texttt{riemannHypothesis},
\texttt{LFunction\_ne\_zero\_of\_half\_lt\_re}, and
\texttt{LFunction\_eq\_zero\_re} in \texttt{ALPOGE\_Main.lean} have the
definitions printed in \texttt{ALPOGE\_Statements.lean}.  We use these as the
cited input $A_D$, as requested.  The supplied archive contains these exported
declarations and the scan report, but not the imported tree
\texttt{Supplements.Kernels.MainTheoremsAll}; a kernel build of that tree is
not a check performed in this run.  This distinction concerns our evidence,
not the stipulated status of the input theorem.

Every unitary idele-class character of $\mathbb Q$ has the form
$\chi_0|\cdot|^{it}$, where $\chi_0$ is a finite-order Dirichlet character
and $t\in\mathbb R$.  Consequently $A_D$ gives the rank-one Hecke GRH over
$\mathbb Q$ for these characters, not just its finite-order specialization.
Indeed the norm map has a continuous section and compact kernel, which for
$\mathbb Q$ is $\widehat{\mathbb Z}^{\times}$ after the sign convention at
infinity.  A continuous character of this profinite group into the circle has
finite image: its image is a compact quotient of a profinite group, hence
totally disconnected, whereas an infinite closed subgroup of the circle is
the circle.  It therefore factors through some
$(\mathbb Z/N\mathbb Z)^\times$.  The remaining unitary character of
$\mathbb R_{>0}$ is $x^{it}$.  In the convention
$L(s,\chi_0|\cdot|^{it})=L(s+it,\chi_0)$, zeros are translated only in
their imaginary coordinate.  This proves the assertion.  Geometrically,
the finite-order part is a rank-one representation of the arithmetic
fundamental group; the continuous norm twist is an analytic normalization,
not an additional finite-monodromy motive.  The result does not cover
arbitrary characters over a nontrivial number field.

\subsection*{MT-M01. A nonzero counterexample to the draft's general zero-count lemma}
Masullo, Lemma 9.3.2.4, PDF p.~425, claims a local $O(\log |T|)$ zero count
for the Mellin transform of every $\varphi\in\mathcal S_E(G)$.
We checked the image of the original PDF page, including the displayed
Jensen inequality and its last inference.  The quantifier is present in the
PDF; this is not an inference from the model transcription.

Take $G=\mathbb R_{>0}$, trivial character, Haar measure $dt/t$, and set
\[
 g_a(u)=\frac{1}{\sqrt{4\pi a}}\exp\left(-\frac{u^2}{4a}\right),
 \qquad a=1+i,\qquad
 \varphi(e^u)=\frac{g_{1+i}(u)-g_{1-i}(u)}{2i}.
\]
The square roots use the branch with positive real part.  This is a real,
nonzero smooth function.  For every real $b$ and every nonnegative integer
$j$, the $j$th derivative of $e^{bu}g_{1\pm i}(u)$ is a polynomial in $u$
times $e^{bu-u^2/(4(1\pm i))}$.  Its absolute value is bounded by a
polynomial times $e^{bu-u^2/8}$, hence is rapidly decreasing.  Thus
$t^b\varphi(t)$ is Schwartz in the logarithmic coordinate for every $b$;
this is exactly the condition of Definition D.3.3.1, PDF p.~804 (also
\S1.1.2, PDF p.~11).  In particular $\varphi\in\mathcal S_E(G)$.

For $\operatorname{Re}a>0$, the Gaussian integral and analytic continuation
in $a$ give
\[
 \int_{\mathbb R}g_a(u)e^{su}\,du=e^{as^2},\qquad s\in\mathbb C.
\]
For completeness, the identity first follows by completing the square for
positive real $a$ and real $s$; both sides are entire in $s$ by Gaussian
domination on compact sets, and holomorphic in $a$ on
$\operatorname{Re}a>0$, where $\operatorname{Re}(1/a)>0$.  The identity
theorem proves the claimed continuation.  Therefore
\[
 \mathcal M\varphi(s)=e^{s^2}\sin(s^2).
\]
Its zeros on the positive imaginary axis are
$i\sqrt{k\pi}$, $k=1,2,\ldots$, and are simple.  For $I=[-1,1]$ and
$T\ge2$, their number in $I+i[T,T+1]$ is
\[
 \#\{k:T^2/\pi\le k\le(T+1)^2/\pi\}
   =\frac{2T+1}{\pi}+O(1).
\]
This is not $O(\log T)$.  The function is nonzero, so removing the zero
function from the lemma does not repair it.  The proof on p.~425 needs a
\emph{lower} bound for the modulus at a Jensen centre, whereas rapid decay
supplies upper bounds.  No statement about the zeros of a Hecke L-function
is contradicted by this example.

\subsection*{MT-M02. The displayed approximate identities are proper kernel projections}
Theorem 9.3.2.6, equation (9.3.2.10), PDF p.~426, defines subspaces by
intersecting kernels of all jet evaluations up to height $T$, then claims
their orthogonal projections tend strongly to the identity.  The printed
formula and the asserted inclusion in its proof have both been checked on
the page image.

Here is a counterexample even with a single zero, so the zero-count issue
in MT-M01 is irrelevant.  Put
\[
 g(u)=(4\pi)^{-1/2}e^{-u^2/4},\qquad
 \psi(e^u)=-g'(u).
\]
Then $\psi\in\mathcal S_E(\mathbb R_{>0})$, and integration by parts gives
$\mathcal M\psi(z)=z e^{z^2}$.  Its only zero is $0$, with multiplicity one,
strictly inside the strip for $I=[-1,1]$.  Choose Sobolev index $0$.
The draft's Hilbert norm is then
\[
 \|f\|_H^2=\int_{\mathbb R}|f(e^u)|^2e^{2|u|}\,du.
\]
The functional $\delta_0(f)=\int_{\mathbb R}f(e^u)\,du$ is continuous,
because Cauchy--Schwarz bounds its modulus by
$\|f\|_H(\int e^{-2|u|}\,du)^{1/2}=\|f\|_H$.
It is nonzero, since $\delta_0(g)=1$.  For every $T>0$ the subspace displayed
in (9.3.2.10) is the same proper closed subspace $\ker\delta_0$.
Its orthogonal projector cannot converge to the identity.  In particular,
the Riesz vector $r(e^u)=e^{-2|u|}$ has $\|r\|_H=1$ and is killed by every
one of these projectors.  It belongs to the indicated completion because
$C_c^\infty(\mathbb R)$ is contained in $\mathcal S_E$ in logarithmic
coordinates and is dense in this weighted $L^2$ space: truncate first and
then smooth on compact intervals, where the positive weight is bounded
above and below.  The ranges decrease when more evaluation
conditions are added, whereas the proof on p.~426 uses the opposite
inclusion.

\subsection*{MT-M03. The exact projection repair}
Let $H$ be a Hilbert space, let $(\ell_j)_{j\ge1}$ be continuous linear
functionals, and put
\[
 V_N=\bigcap_{j\le N}\ker\ell_j,\qquad V=\bigcap_{j\ge1}\ker\ell_j.
\]
Then $P_{V_N}\to P_V$ strongly.  The complementary projectors converge
strongly to $P_{V^\perp}$, the projector onto the closed span of the Riesz
vectors of the $\ell_j$.  Thus they give an approximate identity on the
spectral quotient $H/V\simeq V^\perp$.

Here is a full proof.  For fixed $x\in H$, put $y_N=P_{V_N}x$.  For
$M\ge N$, nestedness gives $P_{V_M}y_N=y_M$ and hence
\[
 \|y_N-y_M\|^2=\|y_N\|^2-\|y_M\|^2.
\]
The nonincreasing nonnegative sequence $\|y_N\|^2$ converges, so $(y_N)$ is
Cauchy.  Write its limit as $y$.  For each fixed $N$, all sufficiently late
$y_M$ lie in the closed space $V_N$, so $y\in V$.  Also
$x-y_N\perp V$ for every $N$; passing to the limit gives $x-y\perp V$.
Therefore $y=P_Vx$.  Taking complements proves the second assertion.
Finally the orthogonal complement of the closed span of the Riesz vectors
is precisely $V$, which proves the last identification.  Exhaustion by
height instead of an integer enumeration gives the same proof.

This repairs a functional-analytic direction error.  It does not identify
the common kernel with the closed adelic theta image, prove a Schauder
basis, establish nuclearity of the actual convolution operators, or justify
the geometric-kernel decomposition.  Those are separate statements and
must be checked in the quotient on which they are used.  In Grothendieck's
trace language this is the distinction between an annihilator and its
spectral quotient, before any trace or Fredholm determinant is taken.

\subsection*{MT-M04. An arithmetic replacement for the false general lemma}
For every fixed unitary Hecke character, the number of nontrivial zeros
in a unit-height window is $O_{K,\eta}(\log(|T|+3))$, unconditionally.
The full proof for arbitrary infinity type is NT-Z10 above: polynomial
vertical growth and a uniform lower bound at the Euler-product centre
give Jensen's estimate.  It includes pure norm twists, whose single
translated pole is excluded before applying Jensen at large height.
For a fixed primitive finite-order character a conductor-explicit
primary locator is Greni\'e--Molteni,
\emph{An explicit Chebotarev density theorem under GRH}, Appendix A,
Theorem A.1, author PDF p.~26: that appendix's zero count is unconditional.
Subtract the counts at consecutive heights.  The main term has derivative
$O(\log(|T|+3))$ for a fixed field and character and each remainder has the
same bound.  The arithmetic restriction therefore closes the local
zero-count gate, including infinite-order characters, without using GRH.
It supplies a route around MT-M01 for arithmetic divisors, but does not
establish Lemma 9.3.2.5 for arbitrary Schwartz convolution ideals.  The
constant in the infinite-order statement is fixed-character, not a
claimed uniform archimedean-conductor constant.

\subsection*{MT-ART01. A nonmonomial GL\texorpdfstring{$_2$}{2} consequence: odd Artin forms}
Assume $M_H$.  Let $\rho:G_{\mathbb Q}\to\operatorname{GL}_2(\mathbb C)$
be continuous, irreducible, odd and of finite image.  Then its Artin
L-function is entire and all of its nontrivial zeros lie on
$\operatorname{Re}s=1/2$.

The holomorphy input is already a theorem: Khare--Wintenberger,
\emph{Serre's modularity conjecture (I)}, Corollary 10.2(ii), author PDF
p.~21 (with Theorem 10.1(ii), pp.~20--21), associates $\rho$ to a weight-one
cuspidal newform.  Its standard L-function is entire.  Let $L$ be the finite
Galois extension cut out by $\rho$.  Brauer induction writes its character
as an integer linear combination of characters induced from
one-dimensional characters of subgroups of $\operatorname{Gal}(L/\mathbb Q)$.
Artin induction and class field theory therefore express $L(s,\rho)$ as
a finite product of integer powers of finite-order Hecke L-functions over
the corresponding fixed fields.  In the open critical strip, $M_H$ makes
each factor holomorphic and nonzero away from the critical line.  The
product, including its negative exponents, is consequently holomorphic and
nonzero there.  Combining this with the independent holomorphy theorem
proves GRH for $L(s,\rho)$.

This includes projective image $A_5$: no monomial expression with only
nonnegative exponents is required.  It is a new transfer in this catalogue,
not a new proof of the classical modularity theorem.  In motivic terms it
concerns odd rank-two Artin motives and their weight-one automorphic
realizations.  It gives no GRH for the weight-two L-function of an arbitrary
non-CM elliptic curve.  Nor does it prove holomorphy of arbitrary even
icosahedral or higher-dimensional Artin representations.  The inherited
Brauer divisor-support statements in RETURN\_ZETA\_TWO are the general
input behind this specialization.

\subsection*{MT-ART02. The missing positivity is a finite matrix at each zero}
Let $L/K$ be a finite Galois extension with group $G$, and let
$0<\operatorname{Re}s_0<1$.  By Brauer meromorphy define
\[
 n_\chi=\operatorname{ord}_{s=s_0}L(s,\chi,L/K)\in\mathbb Z,
 \qquad\Theta_{s_0}=\sum_{\chi\in\operatorname{Irr}(G)}n_\chi\chi.
\]
Form the $|G|$ by $|G|$ matrix
\[
 K_{g,h}=\Theta_{s_0}(g h^{-1}).
\]
Then all Artin L-functions for $L/K$ are holomorphic at $s_0$ if and only
if $K$ is positive semidefinite.  The matrix is Hermitian because
$\chi(g^{-1})=\overline{\chi(g)}$ and $n_\chi$ is real.

To prove the equivalence, consider the central convolution operator
$T=|G|^{-1}K$ on $\mathbb C[G]$.  On the matrix-coefficient space of an
irreducible character $\chi$, Schur orthogonality makes $T$ scalar with
eigenvalue $n_\chi/\chi(1)$ (interchanging $\chi$ and its dual under the
opposite convolution convention only permutes the list).  Each such space
has dimension $\chi(1)^2$ and together they exhaust the regular
representation.  Thus $K\succeq0$ is equivalent to $n_\chi\ge0$ for every
$\chi$, which is exactly absence of a pole of every factor.

The restrictions of $\Theta_{s_0}$ to cyclic subgroups are already actual
characters.  Indeed, for a cyclic subgroup $C$ and a one-dimensional
character $\eta$ of $C$, Frobenius reciprocity and Artin induction give
\[
 \langle\Theta_{s_0}|_C,\eta\rangle_C
   =\operatorname{ord}_{s=s_0}L(s,\eta,L/L^C)\ge0,
\]
because the right-hand side is an order of a holomorphic Hecke function in
the strip.  This proves positivity on every such restriction; it does not
prove the full matrix is positive.  Also
\[
 \Theta_{s_0}(1)=\sum_\chi n_\chi\chi(1)
   =\operatorname{ord}_{s=s_0}\zeta_L(s).
\]
The point of the criterion is the exact remaining object: a central
positive kernel in the Artin character ring.  $M_H$ removes all tests off
the critical line, while neither zero location nor cyclic positivity
supplies the full critical-line kernel.  An Alp\"oge-type positivity
argument would have to produce this kernel or an equivalent positive
decomposition, with the actual Artin orders as coefficients.

\subsection*{MT-ART04. All induced-linear positivity tests still leave a gap}
There is a virtual character which passes every positivity inequality
obtained by inducing a one-dimensional character of a subgroup, but is not
an actual character.  This refutes an abstract reconstruction rule; we do
not claim that this virtual character occurs as the orders of actual
Artin L-functions at a point.

Let $G=A_5$, and let $\chi$ be its four-dimensional augmentation character
for the natural action on five letters.  Define
\[
             \Theta=\operatorname{reg}_G-5\chi.
\]
The natural permutation action is doubly transitive, so its permutation
character has inner-product norm squared two.  Its invariant subspace is
one-dimensional; the orthogonal complement therefore has character
$\chi$ with norm squared one, proving that $\chi$ is irreducible.
The coefficient of $\chi$ in $\Theta$ is $4-5=-1$, so $\Theta$ is not an
actual character and its full matrix from MT-ART02 is not positive.

Nevertheless, for every $H\le G$ and every linear character $\eta$ of $H$,
\[
       \left\langle\Theta,\operatorname{Ind}_H^G\eta\right\rangle_G
                    \ge0.                                      \tag{MT.1}
\]
Here are all the details.  The conjugacy classes of $A_5$ have sizes
$1,15,20,12,12$, as follows by counting the double transpositions,
three-cycles and the two five-cycle classes.  A normal subgroup is a union
of these classes containing the identity, and its size divides $60$.
No proper nontrivial sum of these class sizes containing $1$ divides $60$;
thus $A_5$ is simple.  An action on cosets of a proper subgroup of index
$2,3,$ or $4$ would be an injective homomorphism into the corresponding
symmetric group, which is impossible by order.  Lagrange's theorem leaves
only the possible proper indices
\[
                    k\in\{5,6,10,12,15,20,30,60\}.
\]
Put $m=\langle\chi|_H,\eta\rangle_H$.  Since $\chi|_H$ has dimension four,
$0\le m\le4$.  Frobenius reciprocity says that
$\operatorname{Ind}_H^G\eta$, of dimension $k$, contains $m$ copies of
$\chi$, so $4m\le k$.  These two bounds imply $k-5m\ge0$ for every index
in the list except possibly $k=12$.  At $k=12$, $|H|=5$, and $H$ acts
regularly on the five letters.  The restriction of $\chi$ is the sum of
the four distinct nontrivial characters of $C_5$, so $m\le1$ and again
$k-5m\ge0$.  Finally
\[
 \left\langle\operatorname{reg}_G,
       \operatorname{Ind}_H^G\eta\right\rangle=k,
\]
proving (MT.1) for proper $H$.  For $H=G$, simplicity and nonabelianity
make the trivial character the only linear character, and its coefficient
in $\Theta$ is one.  This proves (MT.1) in all cases.

For comparison with the methods, Artin induction and Hecke holomorphy
give exactly inequalities of the form (MT.1) for an actual local Artin
order character.  The example proves that these inequalities alone do not
force a positive character, even before any question of zero location.
Thus a transfer of determinant positivity from the Alp\"oge method would
need additional arithmetic information, not just all the existing
induced-linear nonnegative orders.  The elementary exact computation
\texttt{checks/artin\_cone\_certificate.py} checks the negative matrix
quadratic form: with $v(g)=\chi(g)$ one has $v^*Kv=-900$.

\subsection*{MT-ART03. The classical simple-zero reduction}
Stark, \emph{Some effective cases of the Brauer--Siegel theorem},
Inventiones Mathematicae \textbf{23} (1974), Theorem 3, p.~144, proves
that if $\operatorname{ord}_{s_0}\zeta_L\le1$ and $s_0\ne1$, then all
Artin factors for $L/K$ are holomorphic at $s_0$.  We cite this classical
result rather than count another proof as a finding.  Applied with the
divisor-support theorem under $M_H$, it confines any still-possible pole
inside the strip to a multiple zero of the top-field zeta on the critical
line.  No simplicity assertion for these Dedekind zetas is being made.

\subsection*{MT-Q01. Exact path quadrics in every characteristic}
For $k\ge2$ let
\[
 V_k=\left\{(x_1,\ldots,x_k):\sum_{i=1}^{k-1}x_ix_{i+1}=0\right\}
       \subset\mathbb A^k_{\mathbb Z},\qquad
 r=\lfloor k/2\rfloor,\quad d=\lceil k/2\rceil.
\]
Then for every finite field $\mathbb F_q$, including characteristic two,
\[
 \#V_k(\mathbb F_q)=q^{k-1}+(q-1)q^{d-1}.                 \tag{MT.2}
\]
In the relative Grothendieck ring of varieties over $\mathbb Z$ one has
\[
                 [V_k]=\mathbb L^{k-1}+\mathbb L^d-\mathbb L^{d-1}.
                                                               \tag{MT.3}
\]
Here $\mathbb L=[\mathbb A^1_{\mathbb Z}]$; the identity is an identity of
classes, not a claim of a direct-sum splitting of all integral realizations.

To prove these formulas, group the odd coordinates into $x\in\mathbb A^d$
and the even coordinates into $y\in\mathbb A^r$.  The quadratic form
becomes $x^{\mathsf T}Ay$, where the integral $d$ by $r$ matrix $A$ has an
$r$ by $r$ triangular minor with diagonal entries one.  It has a left
inverse over $\mathbb Z$ and its columns extend to a unimodular basis of
$\mathbb Z^d$.  Thus after an invertible integral change of the $x$
coordinates the equation is $\sum_{j=1}^r x_jy_j=0$, with $d-r$ free
$x$ coordinates.  Over $y=0$ its fibre is $\mathbb A^d$.  Over
$\mathbb A^r\setminus\{0\}$, the map of vector bundles
$\mathcal O^d\to\mathcal O$ given by $(y_1,\ldots,y_r,0,\ldots,0)$ is
surjective, and its kernel is a vector bundle of rank $d-1$.
It is Zariski locally trivial on the opens where a coordinate $y_j$ is
invertible.  Additivity and the vector-bundle relation give
\[
 [V_k]=\mathbb L^d+(\mathbb L^r-1)\mathbb L^{d-1},
\]
which is (MT.3).  Counting the same fibres over $\mathbb F_q$ proves
(MT.2) directly, without a purity theorem.

Consequently the exact error term in the unweighted orthogonality identity
of the supplied Fable sketch is
\[
             q\,\#V_k(\mathbb F_q)=q^k+(q-1)q^d.
\]
For a fixed nonzero additive-character parameter its quadratic sum is
$q^d$, by summing first over $x$.  This gives square-root size when $k$ is
even; for odd $k$ the extra factor $q^{1/2}$ reflects the one-dimensional
radical.  It explains the sign and the exact constant for this family.
The general weighted operator estimate in
\texttt{OURS\_FABLE\_PROGRAM\_SUGGESTIONS\_20261120\_EXPLAINED.md},
\S3, remains a separate inherited result.  This calculation does not prove
the sign or norm assertions for arbitrary sheaves in Alp\"oge's steps 3
and 5, nor identify this elementary model with all the constructions in
the omitted Lean dependency tree.

\subsection*{MT-Q02. Its arithmetic zeta function and shifted zero lines}
The exact all-prime formula in MT-Q01 gives
\[
 Z(V_k/\mathbb F_q,T)
     =\frac{1-q^{d-1}T}{(1-q^{k-1}T)(1-q^dT)},
 \qquad
 \zeta_{V_k}(z)
     =\frac{\zeta(z-k+1)\zeta(z-d)}{\zeta(z-d+1)}.       \tag{MT.4}
\]
Indeed apply (MT.2) over every $\mathbb F_{q^m}$ in
$Z=\exp(\sum_{m\ge1}\#V_k(\mathbb F_{q^m})T^m/m)$ and use
$-\log(1-aT)=\sum_{m\ge1}a^mT^m/m$.  Multiplying the resulting factors
over rational primes proves the second identity first in a right
half-plane of absolute convergence, and then meromorphically.
Under $A_D$ the nontrivial divisor of this function is supported on
\[
    \operatorname{Re}z=k-\tfrac12,\qquad
    \operatorname{Re}z=d+\tfrac12,\qquad
    \operatorname{Re}z=d-\tfrac12.
\]
Divisor orders are signed and cancellations are kept exactly as in
(MT.4); the denominator's zeros are possible poles, not additional zeros.
The shifted pole of $\zeta$ at $1$ and shifted trivial zeros are recorded
by the same explicit formula, separately from this nontrivial-divisor
assertion.  This is an all-characteristic, concrete Artin--Tate example
linking the path model of the method with arithmetic zeta functions and
the Grothendieck ring.  It illustrates why a mixed-weight scheme has
several shifted lines rather than a single universal critical line.

\subsection*{MT-L01. Damped Euler products and Adams operations}
Let $\rho$ be a finite-image complex representation of $G_K$ with no
invariant vectors, of dimension $d$.  At each finite place $v$ put
$A_v=\rho(\operatorname{Frob}_v)|_{V^{I_v}}$.  Its eigenvalues have
absolute value one.  For $|c|\le1$ define, initially for $\operatorname{Re}s>1$,
\[
 D_c(s,\rho)=\prod_v\det(1-cA_v(Nv)^{-s})^{-1}.
\]
If $L(s,\rho)$ is holomorphic and nonzero on
$H=\{\operatorname{Re}s>1/2\}$, then $D_c$ has a holomorphic nonvanishing
continuation to $H$, given by
\begin{align*}
 D_c(s,\rho)&=\exp\{c\log L(s,\rho)+E_c(s)\},\\
 E_c(s)&=\sum_v\sum_{k\ge2}
   \frac{c^k-c}{k}\operatorname{tr}(A_v^k)(Nv)^{-ks}.
\end{align*}
The branch of $\log L$ is fixed by the Euler series in
$\operatorname{Re}s>1$.

For a full proof, expand $-\log\det(1-X)$ in powers of $X$ in the initial
half-plane.  The terms with $k=1$ cancel in
$\log D_c-c\log L$, giving the displayed $E_c$.  On a compact subset with
$\operatorname{Re}s\ge1/2+\varepsilon$, its terms in absolute value sum to
at most
\[
 2d\sum_v\sum_{k\ge2}(Nv)^{-k(1/2+\varepsilon)}
 \le \frac{2d}{1-2^{-1/2-\varepsilon}}
       \sum_v(Nv)^{-1-2\varepsilon}<\infty.
\]
Normal convergence, also uniform in $|c|\le1$, proves that $E_c$ is
holomorphic in $s$.  A nonvanishing holomorphic function on the simply
connected half-plane $H$ has a holomorphic logarithm, uniquely normalized
by its existing branch on the right.  The exponential above is therefore
holomorphic, nonzero, and agrees with the product where it converges.
This proves the continuation.

Under $M_H$, Brauer induction supplies the needed holomorphic nonvanishing
on $H$ for $\rho$ with no invariants: all Hecke factors have no zeros
there, and their orders at $1$ combine to
$-\dim V^{G_K}=0$.  This does not require Artin holomorphy on the boundary
line.  Under $A_D$ the same argument applies to sums of nonprincipal
Dirichlet characters over $\mathbb Q$.

The traces $\operatorname{tr}(A_v^k)$ are the local Adams-power characters;
the proof separates the first prime term from all higher Adams terms,
whose series converge in the larger half-plane.  This is a concrete
$\lambda$-ring link, rather than a claim that positivity of coefficients
follows from the formal operations.  For $c=1/m$ the phenomenon is closely
related to Gun--Hazra--Sahu, \emph{On holomorphy and non-vanishing of Artin
L-functions} (2025), Theorem 5.  The displayed formula covers all complex
$|c|\le1$ and makes its hypotheses explicit.

\subsection*{MT-L02. One irrational parameter detects the zero-free half-plane}
Fix an irrational real number $0<c<1$ and a nontrivial irreducible Artin
representation $\rho$.  Then the following are equivalent:
\begin{enumerate}
 \item $L(s,\rho)$ is holomorphic and nonzero on $H$;
 \item the Euler product $D_c(s,\rho)$ admits a single-valued meromorphic
       continuation to $H$.
\end{enumerate}
In the second assertion meromorphic is a global, single-valued condition;
a branch on a slit domain is not enough.

MT-L01 proves the forward implication.  Conversely $L$ is meromorphic by
Brauer induction and $E_c$ is holomorphic on $H$ by the same normally
convergent series without any hypothesis about the zeros of $L$.
Logarithmic differentiation in $\operatorname{Re}s>1$ gives
\[
 \frac{D_c'}{D_c}=c\frac{L'}L+E_c'.
\]
Both sides are meromorphic on connected $H$ if the asserted continuation
of $D_c$ exists, so the identity holds on $H$.  Taking residues at any
$s_0\in H$ gives
\[
 \operatorname{ord}_{s_0}D_c
      =c\operatorname{ord}_{s_0}L.
\]
Both orders are integers, and $c$ is irrational.  Both must be zero.
Thus $L$ is holomorphic and nonvanishing on $H$, proving the equivalence.
If instead $c=a/b$ in lowest terms, the same equality proves
$b\mid\operatorname{ord}_{s_0}L$ everywhere in $H$.

GRH in the meromorphic-divisor sense for the nontrivial irreducible Artin
factors is therefore equivalent, by the functional equation, to these
irrational-parameter continuation tests.  To include the trivial factor one
must separately add RH for $\zeta_K$; that input is already part of $M_H$.
This is an analytic multiplicity test.  It neither constructs the
continuation from arithmetic lattices nor removes poles on the critical
line.  With a trivial summand of multiplicity $r$, the known pole at $1$
contributes local monodromy $e^{-2\pi i c r}$; this is why the no-invariants
hypothesis was retained rather than suppressed.

\subsection*{MT-ART05. The exact induced-linear cone for \texorpdfstring{$A_5$}{A5}}
The coarse counterexample MT-ART04 can be sharpened to a degree-three
virtual character. More usefully, the same calculation gives explicit
lower bounds for any pole of an actual $A_5$ Artin function. These are
group-theoretic deductions, not realizations of the virtual examples as
arithmetic divisor orders.

Write the irreducible characters in the order
$1,\chi_3,\chi'_3,\chi_4,\chi_5$, and put
\[
 \Theta=a+b\chi_3+c\chi'_3+d\chi_4+e\chi_5,
 \qquad a,b,c,d,e\in\mathbb Z.
\]
Then the inequalities
$\langle\Theta,\operatorname{Ind}_H^{A_5}\eta\rangle\ge0$ for
every subgroup $H$ and every linear character $\eta$ are equivalent to
the following seven inequalities:
\[
 a\ge0,\quad e\ge0,\quad a+d\ge0,\quad b+c\ge0,\quad
 b+d+e\ge0,\quad c+d+e\ge0,\quad b+c+d\ge0.       \tag{MT-A5-cone}
\]
If $r=\Theta(1)=a+3b+3c+4d+5e$, then $d<0$ forces $r\ge3$,
and $b<0$ or $c<0$ forces $r\ge4$. These thresholds are sharp for
this abstract cone: respectively
\[
       1+\chi_3+\chi'_3-\chi_4,
       \qquad -\chi_3+\chi'_3+\chi_4               \tag{MT-A5-ex}
\]
have degrees three and four and satisfy every inequality.

\begin{proof}
We give the character and subgroup details so the finite assertion has
no unreported character-table oracle. The five conjugacy classes have
orders $1,2,3,5,5$ and sizes $1,15,20,12,12$. The natural four-dimensional
augmentation character has values $(4,0,1,-1,-1)$ and norm one, as in
MT-ART04. Its real representation is oriented because all permutations
are even. The Hodge-star decomposition of its exterior square has two
invariant real three-dimensional spaces. The exterior-square character
has values $(6,-2,0,1,1)$ by
\[
 \chi_{\wedge^2 V}(g)=\tfrac12(\chi_V(g)^2-\chi_V(g^2)).
\]
Its norm squared is two and it has no invariant vectors, so the two
three-dimensional summands are distinct irreducibles. Simplicity of
$A_5$ makes each faithful and its determinant trivial. An order-two or
order-three element in either corresponding $\mathrm{SO}_3$ representation
has trace $-1$ or $0$. At order five the two traces are
$\varphi=(1+\sqrt5)/2$ and $\bar\varphi=(1-\sqrt5)/2$; their sum is
one and each is the trace of a nontrivial order-five rotation.

There are six Sylow-five subgroups. Their normalizers have order ten
and are dihedral, since the centralizer of a five-cycle is its own
cyclic subgroup. The permutation action on the six subgroups has
values $(6,2,0,1,1)$: an order-five element fixes its unique Sylow
subgroup, an order-three element fixes none, and double counting gives
$6\cdot5/15=2$ fixed subgroups for an involution. Removing the constants
gives an irreducible character $\chi_5=(5,1,-1,0,0)$, of norm one.
The squared dimensions $1+9+9+16+25=60$ exhaust the character table.

Every proper subgroup is contained in an $A_4$, a $D_5$ of order ten,
or an $S_3$ of order six. Here is an elementary verification. Simplicity
excludes proper indices two, three and four. The remaining possible
proper subgroup orders divide sixty and belong to
$\{1,2,3,4,5,6,10,12\}$. There are no elements of order four, six or
ten. Thus subgroups of orders four, six and ten have types
$V_4,S_3,D_5$. The centralizer of an involution is its unique $V_4$,
whose normalizer is $A_4$. In a subgroup of order twelve, a normal
Sylow-three subgroup is impossible since the ambient normalizer of
$C_3$ has order six. There are therefore four Sylow-three subgroups,
contributing eight nonidentity elements. Three Sylow-two subgroups
would contribute nine further involutions, because distinct $V_4$'s
have disjoint nonidentity sets. Hence the Sylow-two subgroup is normal
and the order-twelve subgroup is its $A_4$ normalizer. The smaller
cyclic and Klein subgroups are contained in these listed groups.

Let $u,\bar u$ be the nontrivial linear characters of $A_4$,
$v$ its irreducible character of degree three, and let $\epsilon$
denote a sign character of either $D_5$ or $S_3$. The restriction
decompositions obtained from the displayed character values are
\[
\begin{array}{c|ccc}
 & A_4 & D_5 & S_3\\\hline
1&1&1&1\\
\chi_3&v&\epsilon+v_1&\epsilon+v_2\\
\chi'_3&v&\epsilon+v'_1&\epsilon+v_2\\
\chi_4&1+v&v_1+v'_1&1+\epsilon+v_2\\
\chi_5&u+\bar u+v&1+v_1+v'_1&1+2v_2
\end{array}                                                   \tag{MT-A5-res}
\]
Here $v_1,v'_1$ are the two dihedral characters of degree two and
$v_2$ is the standard $S_3$ character. The nonlinear characters in
these three groups are themselves induced from linear characters:
$v$ from a nontrivial character of $V_4$, $v_1,v'_1$ from $C_5$,
and $v_2$ from $C_3$. Transitivity of induction and Frobenius reciprocity
therefore show that all induced-linear tests imply nonnegative
multiplicities in every column of (MT-A5-res). Together with the
global trivial-character test $a\ge0$, these are exactly the seven
inequalities (MT-A5-cone), after deleting redundant inequalities.
Conversely those seven inequalities make every column an actual
character. Restricting further gives actual characters on all proper
subgroups, hence every required linear inner product is nonnegative.
The global linear test is again $a\ge0$. This proves the equivalence.

For the degree bounds, first put $d=-k<0$. Then $a\ge k$,
$b+c\ge k$, and $b,c\ge k-e$. If $e=0$, these imply
$r\ge k+6k-4k=3k\ge3$. If $e\ge1$, the first two inequalities give
$r\ge k+3k-4k+5e\ge5$. Next put $b=-k<0$. If $d\ge0$, we have
$c\ge k$, $d+e\ge k$, whence
$r\ge4d+5e\ge4k$. If $d=-j<0$, then
$a\ge j$, $c\ge k+j$, and $e\ge k+j$, giving
$r\ge j+3j-4j+5(k+j)\ge4k$. The argument for $c<0$ is identical.
Finally direct substitution verifies both examples (MT-A5-ex).
\end{proof}

\subsection*{MT-ART06. Pole multiplicity gates for actual \texorpdfstring{$A_5$}{A5} extensions}
Let $L/K$ be an $A_5$ extension and $s_0\ne1$. Its Heilbronn character
has coefficients equal to the local Artin orders. Hecke holomorphy and
induction give all the tests in MT-ART05, unconditionally. Hence every
Artin factor is holomorphic at a zero of $\zeta_L$ of order at most two.
A pole of the four-dimensional factor requires order at least three
and a zero of $\zeta_K$ at the same point. A pole of either
three-dimensional factor requires order at least four. Conditional on
$M_H$, all remaining nontrivial candidate poles lie on the critical line.

At the central point the lower bounds strengthen to six and eight,
respectively. To prove this final refinement, the permutation actions
of degrees five and six give
\[
 L(s,\chi_4)=\frac{\zeta_{L^{A_4}}(s)}{\zeta_K(s)},\qquad
 L(s,\chi_5)=\frac{\zeta_{L^{D_5}}(s)}{\zeta_K(s)}.
\]
The completed Dedekind zeta functional equation has sign $+1$, and
its completion factors are finite and nonzero at $1/2$. Consequently
the central orders $a,d,e$ are even integers, with $a,e\ge0$ by
(MT-A5-cone). If $d=-k<0$, then $k\ge2$; the preceding proof gives
$r\ge3k\ge6$ when $e=0$ and $r\ge5e\ge10$ otherwise.
If $b=-k<0$ and $d\ge0$, the nonnegative even integers $d,e$
satisfy $d+e\ge k\ge1$, so $r\ge4d+5e\ge8$.
If $d=-j<0$, then $j\ge2$ and the preceding bound
$r\ge5(k+j)$ is stronger. The same holds for $c<0$.
This proves all assertions; no assumption on simplicity of zeros has
been introduced. The degree-three and degree-four virtual examples
only show the limits of the indicated character tests away from the
additional central parity constraints.

These refinements put the residual Artin question into explicit local
inequalities on a finite Galois quotient, within the SGA 1 V \S7
framework. They sharpen ORDER item 42 and ZETA\_TWO ZT02 without
claiming general Artin holomorphy or constructing actual negative
Artin divisor orders.

\subsection*{MT-H01. A unitary Hecke family outside the algebraic-CM normalization}
Let $K=\mathbb Q(\sqrt5)$, $\epsilon=(1+\sqrt5)/2$, and let $\sigma$
be its nontrivial automorphism. Put $t_0=\pi/\log\epsilon$.
For each nonzero integer $k$ there is a conductor-one unitary Hecke
character defined on fractional ideals by
\[
       \eta_k((\alpha))=
       \left|\frac{\alpha}{\sigma(\alpha)}\right|^{ik t_0}.
                                                               \tag{MT-H1}
\]
It has infinite order, $\eta_k^\sigma=\eta_{-k}$, and is trivial
on the rational ideles. No positive power of $\eta_k$ becomes an
algebraic Hecke character after multiplying by a complex norm power
and a finite-order character. Thus the full unitary-Hecke input $M_H$
has an explicit scope beyond the algebraic Hecke characters used in
the potentially CM motive transfer.

\begin{proof}
The ring of integers is $\mathbb Z[\epsilon]$, with discriminant five.
Minkowski's ideal-class bound is $\sqrt5/2<2$, so every fractional
ideal is principal. An exact classical locator for the bound is
Milne, \emph{Algebraic Number Theory}, Theorem 4.3 (v3.08), PDF p.~72.
The units are $\{\pm\epsilon^j:j\in\mathbb Z\}$. For a direct
verification, a positive unit $u$ with $1<u<\epsilon$ would have
integer trace $u+u^{-1}$ in $(2,\sqrt5)$ if its norm were one, or
$u-u^{-1}$ in $(0,1)$ if its norm were minus one; neither interval
contains an integer. Multiplying an arbitrary positive unit by a
power of $\epsilon$ reduces it to $[1,\epsilon)$, so it is a power
of $\epsilon$, and signs give the assertion.

Multiplying $\alpha$ by a unit changes the logarithm of the absolute
ratio in (MT-H1) by $2j\log\epsilon$. Its phase therefore changes by
$2\pi i k j$, so (MT-H1) is well-defined and multiplicative. It is
unitary. The classical ideal/infinity-type description of Hecke
characters makes it an unramified finite-place character, with
archimedean exponents $\pm ikt_0$ (their simultaneous sign depends
on the ideal-to-idele convention). The unit compatibility just checked
is exactly what is required for this description. Rational ideles
have equal archimedean components and equal valuations at conjugate
places; their images are trivial. For a ramified rational place this
also follows here from the generator $\sqrt5$, whose absolute ratio
is one. Conjugating (MT-H1) inverts the ratio.

On the positive archimedean subgroup the exponents of any algebraic
Hecke character are integers, up to an overall convention sign.
A norm power adds the same complex number to the two exponents;
a finite-order character is trivial on this connected subgroup.
If a positive power $m$ of (MT-H1) had the asserted algebraic
normalization, subtraction of its two exponents would give
$2mikt_0$ equal to the difference of two integers. This is impossible.
The same nonconstant archimedean restriction proves infinite order.
\end{proof}

This does not assert a transcendence theorem for individual Hecke
values, or rule out every conceivable generalized motive. It precisely
separates the character class covered by CM01 from the broader unitary
character class quantified in Masullo's endpoint. A proposed transport
of the CM theta-value and algebraic-height steps must therefore address
this extra archimedean data; norm twisting alone cannot do so.

\subsection*{MT-H02. An explicit nonholomorphic \texorpdfstring{$\mathrm{GL}_2$}{GL2} family}
For $k\ge1$, the classical Hecke--Maa\ss\ correspondence associates to
$\eta_k$ a cuspidal dihedral Maa\ss\ newform $g_k$ of weight zero,
level five and nebentypus $\chi_5$. Its spectral parameter is $kt_0$,
its Laplace eigenvalue is $1/4+k^2t_0^2$, and
\[
       L(s,g_k)=L_K(s,\eta_k).
                                                               \tag{MT-H2}
\]
Consequently $M_H$ gives GRH for every member of this family. Its
entireness is already the classical Hecke theorem. The exact primary
locator for the correspondence, level, spectral parameter and coefficient
identity is Humphries--Khan, \emph{On the Random Wave Conjecture for
Dihedral Maa\ss\ Forms}, GAFA 30 (2020), Appendix A.1, printed pp.~106--107 (PDF pp.~73--74), the displayed
formula for $\psi((\alpha))$ and the subsequent formulas for
$t_g$ and $\lambda_{g_\psi}(n)$.

For clarity, the transfer proof is short and complete. MT-H01 gives
a non-norm character: its two archimedean exponents are unequal.
The cited correspondence is therefore cuspidal, and its formula
$\lambda_{g_k}(n)=\sum_{N\mathfrak a=n}\eta_k(\mathfrak a)$ identifies
the absolutely convergent Dirichlet series with (MT-H2), including all
finite primes. Analytic continuation gives the equality everywhere.
The character is unitary, so its nontrivial critical line is
$\operatorname{Re}s=1/2$, which is exactly the assertion supplied by
$M_H$. This proves the stated implication without a generic
non-dihedral Maa\ss\-form hypothesis.

These forms live on the analytic modular curve of level five and in
the automorphic-induction part of the adelic framework. They are an
explicit additional $\mathrm{GL}_2$ branch beside the holomorphic
CM and finite-image weight-one branches, not the missing general
non-CM $\mathrm{GL}_2$ theorem in ORDER item 42.

\subsection*{MT-H03. Exact tensor and symmetric-power products in that family}
Write $F_j(s)=L_K(s,\eta_j)$ for $j\ne0$ and
$F_0(s)=\zeta_K(s)=\zeta(s)L(s,\chi_5)$. Then $F_{-j}=F_j$.
For the actual local two-dimensional parameters $V_j$ of $g_j$, the
following finite Euler-product identities hold:
\[
 L(s,V_m\otimes V_n)=F_{m+n}(s)F_{m-n}(s),\qquad m,n\ge1,
                                                               \tag{MT-H3}
\]
and, for $n\ge0$ and $k\ge1$,
\[
 L(s,\operatorname{Sym}^n V_k)=
 \prod_{j=0}^{\lfloor(n-1)/2\rfloor}F_{(n-2j)k}(s)
 \begin{cases}\zeta(s),&n\text{ even},\\1,&n\text{ odd}.
 \end{cases}                                                   \tag{MT-H4}
\]
The product is empty when $n=0$. Conditional on $M_H$, all nontrivial
zeros of these products lie on $\operatorname{Re}s=1/2$. They are
entire except for a simple pole at one when $m=n$ in (MT-H3), or
when $n$ is even in (MT-H4). This is a covered arbitrarily large
degree family, not a use of an unknown general symmetric-power lift.

\begin{proof}
Conjugating ideals proves $F_{-j}=F_j$. At a split rational prime
the two local parameters of $V_k$ are $z^k,z^{-k}$, where
$z=\eta_1(\mathfrak p)$. The tensor weights are
$z^{m+n},z^{m-n},z^{-m+n},z^{-m-n}$, which gives (MT-H3).
The symmetric-power weights are $z^{(n-2j)k}$ for $0\le j\le n$;
pairing opposite weights leaves one middle weight $1$ precisely for
even $n$, proving (MT-H4) there.

At an inert prime the unique prime ideal is $(p)$ and has character
value one. The two induced Frobenius eigenvalues are $1,-1$ for
every nonzero $k$. Tensoring gives two of each sign. On the
symmetric power, the even-index monomials give eigenvalue $1$ and
the odd-index monomials give $-1$. If $n=2r$, these multiplicities
are $r+1,r$; if $n=2r+1$, they are $r+1,r+1$. These are exactly
the factors on the right sides. If $m=n$, the factor $F_0$ at an
inert prime is $(1-p^{-2s})^{-1}$, so the same check includes that
case.

The only ramified prime is five. The ideal $(\sqrt5)$ has
$\eta_k$-value one and the character is unramified at every finite
place. Thus its local character at that quadratic completion is
trivial, and its induced Weil parameter is $1\oplus\chi_{5,5}$.
The inertia-fixed dimension is one. A tensor product has two
trivial constituents and two $\chi_{5,5}$ constituents; a symmetric
power has respectively $r+1,r$ or $r+1,r+1$ of them. The latter
constituents have no inertia invariants. Hence the ramified Euler
factors also agree exactly, rather than only away from five.
Equivalently the proof is the weight-pair decomposition for diagonal
matrices $\operatorname{diag}(z,z^{-1})$ and an outside coset
acting by interchanging the two basis vectors.

Absolute convergence initially gives the identities for
$\operatorname{Re}s>1$. The right sides give their meromorphic
continuations. All $F_j$, $j\ne0$, are entire Hecke functions with
nonzero boundary value at one; the latter is Ramakrishnan--Valenza,
Theorem 7-28, printed p.~289. The same holds for $L(1,\chi_5)$.
Thus the displayed zeta factors give exactly the asserted simple
poles, with no cancellation at one. Finally $M_H$ locates all the
nontrivial zeros of each unitary Hecke factor. The identities contain
no negative exponents, so they introduce no hidden poles on the
critical line. This proves every assertion.
\end{proof}

The middle summand in (MT-H4) is the trivial character for this
specific family. In a differently normalized dihedral family it
depends on the central character; replacing it universally by either
$1$ or $\chi_K$ would be incorrect. This is the representation-ring
and lambda-operation part of the adelic framework, beside SGA 6 V
\S\S2,7, with the explicit analytic parameters of MT-H01 retained.

\subsection*{MT-ART07. The entire non-real $A_5$ Artin category}
Let $L/\mathbb Q$ be Galois with group $A_5$, and suppose $L$ is not
totally real. Every nontrivial irreducible Artin L-function for this
extension is entire, unconditionally. Under $M_H$ each satisfies GRH.
This closes the non-real $A_5$ specialization of the holomorphy gate;
MT-ART05/06 remain useful for totally real extensions and other base fields.

The holomorphy assertion is classical, not a new finding. An exact primary
locator is Calegari, \emph{The Artin conjecture for some $S_5$-extensions},
\S1, publisher-author PDF p.~2, the paragraph immediately before
Theorem 1.1: it states the stronger automorphy assertion for $A_5$ when
complex conjugation is nontrivial, citing Khare--Wintenberger and Kim's
icosahedral functorial constructions. The two conditions are equivalent
here because a Galois extension has a real embedding exactly when every
embedding is real, equivalently complex conjugation acts trivially.
We cite that classical result rather than repeat its proof.

Here is the complete GRH transfer and its tensor-category form.
For every finite-dimensional representation $V$ of $A_5$, put
$m=\dim V^{A_5}$. Complete reducibility gives
\[
 L(s,V)=\zeta(s)^m
          \prod_{\chi\ne1}L(s,\chi)^{\langle V,\chi\rangle},
 \qquad \langle V,\chi\rangle\in\mathbb Z_{\ge0}.
\]
The product after removing $\zeta^m$ is entire by the cited theorem.
For each irreducible factor, Brauer induction expresses its L-function
as a finite product of integer powers of finite-order Hecke L-functions.
At a point of the open critical strip off its middle line, $M_H$ makes
every such factor finite and nonzero. The product is consequently nonzero
there, including when some Brauer exponents are negative. This proves
the claimed GRH without making a positivity assumption on those exponents.

At $s=1$ every nontrivial irreducible factor is also nonzero. Indeed in a
Brauer expression $\chi=\sum_j a_j\operatorname{Ind}_{H_j}^{A_5}\eta_j$,
the order there is $-\sum_{\eta_j=1}a_j$, by the simple Dedekind pole and
Hecke boundary nonvanishing. Frobenius reciprocity identifies the sum with
$\langle\chi,1\rangle=0$. Thus $L(s,V)$ has a pole of order exactly $m$
at $1$, and has no other poles. Tensor products, duals, exterior and
symmetric powers, and every Schur functor produce representations of the
same finite group, so the same conclusions hold for them after removing
their explicitly determined invariant summands. This is a statement about
the whole Artin tensor category attached to this cover, not about a
general positive-weight motive. The companions in Grothendieck's language
are SGA~1, Expos\'e~V, \S7 (finite monodromy), and SGA~6, Expos\'e~V,
\S\S2,7 (representation-ring operations), together with AG-LAMBDA.

\subsection*{MT-ART08. A finite exceptional set is enough for a holomorphy transfer}
Here is a useful precise comparison lemma. Let $V$ be a self-dual,
finite-image complex Galois representation over a number field $K$.
Suppose an Euler product $F(s)$ converges in a right half-plane and
continues to an entire function, with the same finite local factors as $L(s,V)$ outside a
finite set $S$. Suppose also that, for each $v\in S$, the reciprocal of
the local factor of $F$ is a polynomial $Q_v((Nv)^{-s})$.
Then $L(s,V)$ is entire. No equality of ramified local parameters is
needed for this conclusion.

To prove this, write the Artin factor at $v$ as
$P_v((Nv)^{-s})^{-1}$. Every eigenvalue of Frobenius on the inertia
invariants is a root of unity, since $V$ has finite image. Hence the zeros
of $P_v((Nv)^{-s})$ all lie on $\operatorname{Re}s=0$.
Equality in a right half-plane, followed by meromorphic continuation,
gives
\[
 L(s,V)=F(s)\prod_{v\in S}
             \frac{Q_v((Nv)^{-s})}{P_v((Nv)^{-s})}.
\]
The right side is holomorphic on $\operatorname{Re}s>0$.
The classical completed Artin functional equation is
\[
 A^{s/2}\Gamma_V(s)L(s,V)
   =W A^{(1-s)/2}\Gamma_V(1-s)L(1-s,V),
\]
with $A>0$ and $W\ne0$; the gamma factor is a product of
$\Gamma_{\mathbb R}(s)$, $\Gamma_{\mathbb R}(s+1)$ and
$\Gamma_{\mathbb C}(s)$ with nonnegative integral multiplicities.
An exact locator for these formulas is Prasad--Yogananda,
\emph{A Report on Artin's holomorphy conjecture}, \S3.3,
PDF pp.~8--9, with Brauer induction in Theorem 3.3.
For $\operatorname{Re}s<1$, the numerator
$\Gamma_V(1-s)L(1-s,V)$ is holomorphic, since
$\operatorname{Re}(1-s)>0$. The reciprocal $1/\Gamma_V(s)$ is entire.
The functional equation therefore makes $L(s,V)$ holomorphic on this
second half-plane too. The two half-planes cover $\mathbb C$, proving
the assertion. For a non-self-dual $V$ the identical argument works if
the hypotheses are also supplied for $V^\vee$.

The finiteness of $S$ is an actual hypothesis. Agreement at a set of
places of density one does not give the finite product in the proof.
In particular this lemma cannot by itself upgrade the weak
five-dimensional correspondence in Calegari's Theorem 1.1.
In geometric terms this separates modification at finitely many closed
points of an arithmetic base from a density statement about infinitely
many Frobenius fibres.

\subsection*{MT-S501 and MT-S502. The $S_5$ transfer and its one-value gate}
Let $L/\mathbb Q$ be an $S_5$ extension, with complex conjugation of cycle
type $(12)(34)$. Assume $5$ is unramified and its Frobenius also has this
cycle type. Then the irreducible factors of dimensions $4$ and $6$ are
entire and, under $M_H$, satisfy GRH (MT-S501).

Let $E=L^{A_5}$, let $F\subset L$ have degree $6$ over $\mathbb Q$, and
put $H=EF$. Under $M_H$, the additional single condition
\[
                         \zeta_H(1/2)\ne0
\]
implies that every nontrivial irreducible factor is entire and satisfies
GRH (MT-S502). This is a proved reduction with the displayed extra input,
not a proof of that central nonvanishing.

The classical inputs are Calegari's Theorems 1.1 and 1.2,
publisher-author PDF pp.~2--3; the definition $H=EF$ is explicit in
Theorem 1.2 of arXiv:1112.1152v1, p.~2. The former supplies the
dimension-$4$ and dimension-$6$ holomorphy, and the latter supplies the
remaining case if $\zeta_H$ has no zero in the real interval $(0,1)$.
Under $M_H$ the only possible real zero in that interval is $1/2$.
Thus its interval condition is exactly equivalent to the displayed
single-value condition. Once holomorphy is available, the Brauer
zero-location argument of MT-ART07 proves GRH for every factor.
This is the full proof of both transfers. The one-dimensional factors
are Dirichlet functions and already fall under $A_D$.

The totally real quadratic field $E$ is the base for the odd projective
two-dimensional representations used in the cited theorem. Accordingly
this is a higher-rank consequence of a specific automorphic descent
theorem. It is not obtained by treating every $S_5$ factor as a single
Hecke character. In the Artin-motive programme, the remaining gate in
this family is one central Dedekind value. AG-SIGN controls a nonzero
leading coefficient after a central zero is removed; it does not prove
that the central order is zero.

\subsection*{MT-S503. One unconditional holomorphy example becomes a full GRH example}
Take the splitting field $L$ of
\[
                        x^5-x^3-x^2+x+1.
\]
Calegari, Theorem 5.1, publisher-author PDF p.~15, proves automorphy of
every representation of $\operatorname{Gal}(L/\mathbb Q)=S_5$;
the quintic field has discriminant $1609$. Consequently under $M_H$
every nontrivial irreducible Artin L-function of this explicit field
satisfies GRH, by the same signed Brauer argument. This example uses
that cited theorem, not the local-at-$5$ hypotheses of MT-S501, which
the source explicitly says this polynomial does not satisfy.
The historical bounded-height computation in Theorem 5.3 is part of
the classical proof's evidence; it is not a computation performed in
this run. Geometrically this supplies a concrete finite \(S_5\)-cover
of the arithmetic base whose complete Artin category is covered.

\subsection*{MT-M05. The stated Schauder quotient loses more than its zero evaluations see}
Lemma 9.3.2.5 on original PDF pp.~425--426 asserts a basis for the
quotient by the actual convolution ideal $\varphi*\mathcal S_E(G)$,
with coefficient functionals the Mellin jets at the zeros of
$\mathcal M\varphi$. We inspected the original p.~426 image: there
is no closure over this ideal. The assertion is false even when the
Mellin divisor consists of one simple zero. Thus the counterexample
does not rely on the zero-count defect in MT-M01.

Use $G=\mathbb R_{>0}$, trivial character and $u=\log t$; convolution
becomes ordinary additive convolution. Write
\[
 g_a(u)=(4\pi a)^{-1/2}e^{-u^2/(4a)},\qquad a>0.
\]
Choose $0<b<a$ and set $\varphi(e^u)=-g_a'(u)$,
$h(e^u)=-g_b'(u)$. Both functions belong to the exponential-Schwartz
space, and
\[
 \mathcal M\varphi(s)=s e^{as^2},\qquad
 \mathcal Mh(s)=s e^{bs^2}.
\]
The first function has just the simple zero $0$, lying in the interior
of the strip for $I=[-1,1]$. The only claimed coefficient functional is
$\delta_0$, and $\delta_0(h)=0$.
But $h\notin\varphi*\mathcal S_E$: if $h=\varphi*f$, multiplicativity
of the absolutely convergent Mellin transform gives
$\mathcal Mf(s)=e^{(b-a)s^2}$ for $s\ne0$, and then everywhere by
analyticity. On the imaginary axis this has modulus
$e^{(a-b)\tau^2}$. In contrast, $f(e^u)$ is integrable and therefore
$|\mathcal Mf(i\tau)|\le\int|f(e^u)|\,du$, a contradiction.
Thus the nonzero quotient class of $h$ has every asserted coefficient
zero, which is incompatible with uniqueness in a Schauder expansion.
The still simpler choice $\varphi=g_a$, $h=g_b$ has no Mellin zeros
and a nonzero quotient, so even the empty-divisor case fails.

\subsection*{MT-M06. The closed-ideal repair for every polynomial Gaussian divisor}
In logarithmic coordinates put
\[
 \mathcal E=\{f\in C^\infty(\mathbb R,\mathbb C):
   p_{N,j}(f):=\sup_x e^{N|x|}|f^{(j)}(x)|<\infty
   \text{ for every integer }N,j\ge0\}.
\]
This is $\mathcal S_E(\mathbb R_{>0})$ with its exponential-Schwartz
Fr\'echet topology. Polynomial factors in the usual Schwartz seminorms
are absorbed by increasing $N$. Let
\[
 P(s)=\prod_{\nu=1}^r(s-\rho_\nu)^{m_\nu},\quad
 d=\sum_\nu m_\nu,\qquad
 \varphi=P(-D)g_a,\quad D=d/dx.
\]
Distinct $\rho_\nu\in\mathbb C$, positive integral $m_\nu$, and
$a>0$ are arbitrary; $P=1,d=0$ is allowed. Then
\[
 \overline{\varphi*\mathcal E}
   =J_P:=\{h\in\mathcal E:
          (\mathcal Mh)^{(j)}(\rho_\nu)=0,
          \ 0\le j<m_\nu\},                         \tag{MT-M6}
\]
and the quotient by this \emph{closed} ideal is topologically
$\mathbb C^d$, with exactly the displayed Mellin jets as coordinates.
In particular its coordinate basis is absolute. This repairs the
finite-divisor case after making the closure operation explicit.

\begin{proof}
First, $g_a*\mathcal E$ is dense in $\mathcal E$. We give the full
argument, since density is distinct from the false surjectivity in
MT-M05. Convolution is continuous on $\mathcal E$: the inequality
$e^{N|x|}\le e^{N|x-y|}e^{N|y|}$ bounds each differentiated
convolution by an exponential supremum seminorm times an exponential
$L^1$ seminorm; the latter is bounded by a constant times
$p_{N+1,0}$.

Suppose a continuous linear functional $T$ vanishes on
$g_a*\mathcal E$. For every real $y$ and $\epsilon>0$ the function
$g_{a+\epsilon}(x-y)=g_a*g_\epsilon(\,\cdot-y)$ lies in this range.
It converges to $g_a(x-y)$ in every $p_{N,j}$ as
$\epsilon\downarrow0$, by uniform Gaussian estimates with the
parameter in a compact subinterval of $(0,\infty)$. Hence
$T(g_a(\,\cdot-y))=0$. The function
\[
 B(z)=T\bigl(e^{-x^2/(4a)}e^{zx/(2a)}\bigr)
\]
is entire: on every compact set of $z$, all its proposed complex
derivatives converge in every $p_{N,j}$ by Gaussian domination.
Completing the square gives $B(y)=0$ for real $y$, so the identity
theorem gives $B=0$.

For any Schwartz function $\psi$, Fourier inversion now gives
\[
 T(e^{-x^2/(4a)}\psi(x))
  =\frac1{2\pi}\int_{\mathbb R}\widehat\psi(\xi)
        B(2ai\xi)\,d\xi=0.
\]
This passage under $T$ is justified in the Fr\'echet space itself:
the $p_{N,j}$ seminorm of
$e^{-x^2/(4a)}e^{i\xi x}$ is at most
$C_{N,j,a}(1+|\xi|)^j$, whereas $\widehat\psi$ is rapidly
decreasing. For $h\in C_c^\infty$, take
$\psi=e^{x^2/(4a)}h$, which is again compactly supported and smooth.
Then $T(h)=0$. Finally $C_c^\infty$ is dense in $\mathcal E$:
multiply by a smooth cutoff equal to one on $[-R,R]$ and supported
in $[-2R,2R]$; the product rule and the bound from $p_{N+1,j}$ make
every $p_{N,j}$ of the discarded tail tend to zero. Thus $T=0$.
Hahn--Banach separation proves the asserted density.

Second, the range of $P(-D)$ on $\mathcal E$ is exactly $J_P$.
Integration by parts gives $\mathcal M(-Df)(s)=s\mathcal Mf(s)$,
so inclusion in $J_P$ follows. For the converse, if
$\mathcal Mh(\rho)=0$, define
\[
 f(x)=e^{-\rho x}\int_{-\infty}^x e^{\rho t}h(t)\,dt
     =-e^{-\rho x}\int_x^\infty e^{\rho t}h(t)\,dt.
\]
The equality is precisely the vanishing of the Mellin value.
The first expression controls the negative tail and the second the
positive tail. For any prescribed $N$, using an exponential bound
on $h$ with exponent larger than $N+|\operatorname{Re}\rho|+1$
proves exponential decay of $f$ on both tails; the equation
$f'=h-\rho f$ then gives the same bounds for every derivative.
Therefore $f\in\mathcal E$, and $(D+\rho)f=h$.
Equivalently $-f$ solves $(-D-\rho)(-f)=h$, and its Mellin transform
is $\mathcal Mh(s)/(s-\rho)$. Repeating this division through all
roots, including multiplicities, proves $P(-D)\mathcal E=J_P$.

The Mellin jet functionals are continuous, so $J_P$ is closed.
Since $g_a*\mathcal E$ is dense and $P(-D)$ is continuous,
$P(-D)(g_a*\mathcal E)$ is dense in $P(-D)\mathcal E=J_P$.
Differentiation commutes with convolution, identifying that range
with $\varphi*\mathcal E$. This proves (MT-M6).

For $d=0$ the quotient is zero and no interpolation is required.
For $d>0$, it remains to verify that no coordinate is missing. Given arbitrary
$d$ target jets, ordinary polynomial Hermite interpolation constructs
a polynomial $Q$ of degree below $d$ for which
$Q(s)e^{bs^2}$ has exactly those jets, with any fixed $b>0$.
Indeed multiplication by the nowhere-zero entire function $e^{bs^2}$
is invertible on each finite jet space, so interpolate the resulting
jets of $Q$. The function $Q(-D)g_b$ belongs to $\mathcal E$ and
realizes them. Choosing the targets to be the coordinate vectors
gives a continuous linear section from the finite-dimensional space
$\mathbb C^d$ to $\mathcal E$. The jet map consequently induces
a continuous isomorphism of the quotient with $\mathbb C^d$,
with continuous inverse. Finite-dimensional coordinate expansions are
absolute, completing the proof.
\end{proof}

For $P=s$, MT-M05 therefore exhibits a proper, nonclosed convolution
ideal whose closure is $\ker\delta_0$. The closed quotient really
does have the one coordinate the draft intended. For $P=1$ the
principal ideal is dense but proper, and the closed quotient is zero.
This is a concrete repair rather than a certification of the general
infinite-divisor statement. For arithmetic Hecke divisors the remaining
gate is controlled infinite interpolation and the claimed absolute
basis/nuclear estimates on the correct separated quotient. A local
$O(\log T)$ zero count alone does not prove those estimates.
In cote~1 trace language, Hausdorff completion and the identification
of the spectral quotient precede taking its nuclear trace.

\subsection*{MT-M07. An infinite closed quotient with an absolute Mellin basis}
The closed-ideal repair also admits a complete infinite-divisor example.
Fix $a,\ell>0$ in the logarithmic space $\mathcal E$ of MT-M06, and put
\[
 \varphi(x)=g_a(x-\ell)-g_a(x),\qquad
 \rho_n=\frac{2\pi i n}{\ell}\quad(n\in\mathbb Z).
\]
Then $\mathcal M\varphi(s)=e^{as^2}(e^{\ell s}-1)$, with precisely
the simple zeros $\rho_n$. Let
\[
 (\mathcal P_\ell f)(t)=\sum_{k\in\mathbb Z}f(t+k\ell).
\]
The quotient $\mathcal E/\overline{\varphi*\mathcal E}$ is
topologically $C^\infty(\mathbb R/\ell\mathbb Z)$ under
$\mathcal P_\ell$, and has an absolute Schauder basis whose coefficient
functionals are exactly $f\mapsto\mathcal Mf(\rho_n)$.

\begin{proof}
The periodization series and every derivative converge uniformly on a
fundamental interval, by exponential decay. This gives a continuous map
to the indicated smooth periodic space. Choose
$\chi\in C_c^\infty(\mathbb R)$ nonnegative and positive on
$[-\ell,\ell]$, and set
$\theta=\chi/\mathcal P_\ell\chi$. The denominator is positive,
smooth and periodic, so $\theta$ is smooth, compactly supported and
$\mathcal P_\ell\theta=1$. The map
$R(F)(x)=\theta(x)F(x)$ is a continuous section of periodization,
proving surjectivity.

Let $\Delta_\ell h(x)=h(x-\ell)-h(x)$. Telescoping shows that its
range lies in $\ker\mathcal P_\ell$. Conversely, if
$\mathcal P_\ell f=0$, define
\[
 h(x)=\sum_{k\ge1}f(x+k\ell)
     =-\sum_{k\le0}f(x+k\ell).
\]
These locally normally convergent series define a smooth function.
For $x\ge0$ use the first expression and for $x\le0$ the second.
For every positive integer $N$ and every $j$, summing the geometric
series of bounds $p_{N,j}(f)e^{-N|x+k\ell|}$ proves
$p_{N,j}(h)<\infty$; the $N=0$ case follows from $N=1$.
Thus $h\in\mathcal E$, and direct index shifting gives
$\Delta_\ell h=f$. The range of $\Delta_\ell$ is therefore the
closed space $\ker\mathcal P_\ell$.

By MT-M06, $g_a*\mathcal E$ is dense in $\mathcal E$.
Continuity of $\Delta_\ell$ and commutation with convolution imply
\[
 \overline{\varphi*\mathcal E}
 =\overline{\Delta_\ell(g_a*\mathcal E)}
 =\ker\mathcal P_\ell.
\]
The continuous section $R$ identifies the quotient topology with
the usual smooth periodic topology, rather than merely giving a
vector-space identification.

The Fourier expansion in the convention used here is
\[
 \mathcal P_\ell f(t)=\frac1\ell
       \sum_{n\in\mathbb Z}\mathcal Mf(\rho_n)e^{-\rho_n t}.
\]
Its coefficients are obtained by integrating over a fundamental
interval and summing translates, justified by absolute convergence.
They decrease faster than every power of $|n|$: integrate by parts
arbitrarily many times and use the integrability of every derivative
of $f$. The Fourier series consequently converges absolutely in
every smooth seminorm and equals $\mathcal P_\ell f$ by Fourier
uniqueness. Representatives for the basis are
\[
                     b_n(x)=\ell^{-1}\theta(x)e^{-\rho_n x}.
\]
Indeed $\mathcal Mb_n(\rho_m)=\delta_{m,n}$, by
$\mathcal P_\ell\theta=1$ and integration on one period. Moreover,
$p_{N,j}(b_n)\le C_{N,j}(1+|n|)^j$ since $\theta$ has compact
support. Hence the expansion
$\sum_n\mathcal Mf(\rho_n)b_n$ is absolutely convergent even in
$\mathcal E$, and represents the class of $f$ in the quotient.
The jets determine that class uniquely by the Fourier identification.
Ordering the indices as $0,1,-1,2,-2,\ldots$ gives an absolute
Schauder basis, as asserted.
\end{proof}

The actual ideal still cannot replace its closure: with $0<b<a$,
$h=\Delta_\ell g_b$ has all these Mellin values zero but is not
$\varphi*f$ for any $f\in\mathcal E$, by exactly the Fourier-growth
division argument of MT-M05. Thus this example tests an infinite
divisor with controlled spacing, not just the empty or finite case.

\subsection*{MT-M08. The corresponding trace and positive convolution formula}
For $f\in\mathcal E$, convolution by $f$ descends to the quotient of
MT-M07. On $L^2(\mathbb R/\ell\mathbb Z,dt)$ it is trace class and
\[
 \operatorname{Tr}C_f
   =\sum_{n\in\mathbb Z}\mathcal Mf(\rho_n)
   =\ell\sum_{k\in\mathbb Z}f(k\ell).                 \tag{MT-M8}
\]
All sums converge absolutely. If
$f=h*\widetilde h$ with $h\in\mathcal E$ and
$\widetilde h(x)=\overline{h(-x)}$, this operator is positive and
\[
                  \operatorname{Tr}C_f
                   =\sum_n|\mathcal Mh(\rho_n)|^2\ge0.
\]

To prove the assertions, first use absolute convergence and changes of
variables to periodize a convolution:
\[
 \mathcal P_\ell(f*g)(t)
  =\int_0^\ell(\mathcal P_\ell f)(t-u)
                     (\mathcal P_\ell g)(u)\,du.
\]
This proves descent and identifies the smooth kernel
$K_f(t,u)=(\mathcal P_\ell f)(t-u)$.
The orthonormal Fourier vectors
$\ell^{-1/2}e^{-\rho_n t}$ are eigenvectors with eigenvalues
$\mathcal Mf(\rho_n)$. These eigenvalues decrease faster than every
power, as proved in MT-M07, so their absolute sum is finite.
The diagonal expansion is thus a trace-class expansion and gives
the first trace sum. Evaluating the absolutely convergent Fourier
series of $\mathcal P_\ell f$ at zero gives the second equality
in (MT-M8), including its factor $\ell$. Finally the Mellin transform
of $\widetilde h$ at a purely imaginary argument is the conjugate
of that of $h$; the eigenvalues are therefore the displayed squares.

There is also a direct nuclear realization on the smooth quotient.
Explicitly, put $c_n(F)=\ell^{-1}\int_0^\ell F(t)e^{\rho_nt}\,dt$
and $e_n(t)=e^{-\rho_nt}$. The $c_n$ are an equicontinuous family,
bounded by the uniform norm. In
\[
 C_f=\sum_n(1+|n|)^{-2}\,
       c_n\otimes\bigl((1+|n|)^2\mathcal Mf(\rho_n)e_n\bigr),
\]
the scalars are summable and the vector family is bounded in every
$C^\infty$ seminorm. This is a nuclear decomposition. Fourier partial
sums converge uniformly on bounded sets of $C^\infty$ (apply a
uniform higher-derivative bound to the coefficient tail), so the
finite-rank trace limit is the same absolute sum above.

This supplies a fully checked spectral-quotient, basis and trace
model beside cote~1's nuclear trace language. It does not identify
an arithmetic prime distribution with this lattice or prove the
corresponding basis theorem for the actual Hecke zero set. In that
case the exact infinite interpolation and operator estimates remain
to be supplied on the correct closed quotient.

\subsection*{MT-M09. Even a bounded local zero count does not give the asserted point basis}
Passing to the closed ideal is necessary, but does not prove the
general basis claim. There exists $\varphi\in\mathcal E$ whose Mellin
transform has only simple zeros on the imaginary axis and only
$O(1)$ zeros in every unit-height interval, for which
\[
        Q=\mathcal E/\overline{\varphi*\mathcal E}
\]
has no Schauder basis with the individual zero evaluations as its
coefficient functionals. This holds for every enumeration of the
zeros, and therefore also rules out an absolute Schauder basis of
that specified kind.

\begin{proof}
Put $\epsilon_n=e^{-n^4}$ for $n\ge1$ and define
\[
 H(s)=\prod_{n=1}^{\infty}
       \left(1+\frac{s^2}{n^2}\right)
       \left(1+\frac{s^2}{(n+\epsilon_n)^2}\right),
 \qquad F(s)=e^{as^2}H(s),\quad a>0.
\]
The product converges uniformly on compact sets because the sum of
the absolute deviations of its factors from one does. Away from the
listed zeros the tail logarithms converge locally uniformly, so
there are no additional zeros. Its zeros are exactly
$\pm in,\ \pm i(n+\epsilon_n)$, all simple. Since
$0<\epsilon_n<1/2$, each unit-height interval contains a bounded
number of them.

For $r=|s|$ the product satisfies
\[
 \log|H(s)|\le2\sum_{n\ge1}\log(1+r^2/n^2)\le2\pi r.
\]
For the last bound compare the decreasing function of $x$ with its
integral on $(0,\infty)$; integration by parts gives
$\int_0^\infty\log(1+r^2/x^2)\,dx=\pi r$.
Consequently, on every strip $|\sigma|\le A$,
\[
 |F(\sigma+it)|
 \le \exp\{aA^2-at^2+2\pi\sqrt{A^2+t^2}\}.
\]
Define
\[
 \varphi(x)=\frac1{2\pi}\int_{\mathbb R}F(it)e^{-itx}\,dt.
\]
The Gaussian bound permits every derivative in $x$ under the integral.
Shifting the contour to $\operatorname{Re}s=\sigma$ gives, for
every real $\sigma$ and every $j\ge0$,
\[
 \varphi^{(j)}(x)
   =\frac{e^{-\sigma x}}{2\pi}
      \int_{\mathbb R}(-\sigma-it)^jF(\sigma+it)e^{-itx}\,dt.
\]
The horizontal sides of the contour tend to zero by the displayed
Gaussian estimate, uniformly for $\sigma$ in the finite interval
used in that shift. Choosing $\sigma=N+1$ for $x\ge0$ and
$\sigma=-N-1$ for $x\le0$ proves every seminorm $p_{N,j}$ is finite.
Thus $\varphi\in\mathcal E$. Cauchy estimates applied on slightly
larger strips show that $F(it)$ and all its derivatives are Schwartz.
Fourier inversion gives $\mathcal M\varphi(it)=F(it)$;
both sides are entire, so $\mathcal M\varphi=F$ everywhere.

Every evaluation at a listed zero annihilates the convolution ideal
and its closure, and therefore descends continuously to $Q$.
Let $p=p_{1,0}$ and let
$\overline p([f])=\inf_{j\in\overline{\varphi*\mathcal E}}p(f+j)$
be the corresponding quotient seminorm. For any $f\in\mathcal E$,
\[
 |(\delta_{in}-\delta_{i(n+\epsilon_n)})(f)|
  \le\epsilon_n\int_{\mathbb R}|x||f(x)|\,dx
  \le2\epsilon_n p(f).
\]
The first inequality follows from
$|e^{i\epsilon_n x}-1|\le\epsilon_n|x|$.
Because the difference functional annihilates the closed ideal,
the same bound holds with $p$ replaced by $\overline p$ on $Q$.

Suppose the claimed Schauder basis existed, and denote the basis
vector associated to $\delta_{in}$ by $b_n$. Biorthogonality of
basis vectors and coefficient functionals gives
$\delta_{in}(b_n)=1$ and
$\delta_{i(n+\epsilon_n)}(b_n)=0$. Hence
\[
                 \overline p(b_n)\ge(2\epsilon_n)^{-1}.
\]
Take the single test function $g_b\in\mathcal E$ for any fixed
$b>0$. Its coefficient at $\delta_{in}$ is $e^{-bn^2}$.
The corresponding terms of its proposed basis expansion satisfy
\[
 \overline p\bigl(e^{-bn^2}b_n\bigr)
        \ge\frac12 e^{n^4-bn^2}\longrightarrow\infty.
\]
Terms of a convergent series in a locally convex space must tend
to zero. In every enumeration these infinitely many distinct terms
occur arbitrarily far along the series, giving a contradiction.
This proves the assertion.
\end{proof}

This is a counterexample for general exponential-Schwartz Mellin
divisors, not a claim that such clustered zeros occur for an actual
Hecke L-function. It pinpoints why NT-Z10/MT-M04 cannot by themselves
establish the arithmetic Schauder step: even an $O(1)$ local count
does not control the norm of individual interpolation vectors.
The next possible repair must incorporate quantitative separation,
or group nearby evaluations into divided-difference blocks and prove
convergence of the resulting block expansion. Neither property is
supplied by a zero-location hypothesis alone in the present argument.
